%% Font size %%
\documentclass[11pt]{article}

%% Load the custom package
\usepackage{Mathdoc}

%% Numéro de séquence %% Titre de la séquence %%
\renewcommand{\centerhead}{}

%% Spacing commands %%
\renewcommand{\centerhead}{%
\textbf{Chapitre 13 -- Additions décimales -- DS 3 : Problèmes}
}
\begin{document}


\begin{center}
\duree{55 minutes} 
\coefficient{2}
\calculatrice{0}
\brouillon
\soin
\recherche 
\copieseparee{0}
\end{center}

\phantom{0}
\vspace{.5cm}

\entetedevoirs{20}

\begin{exercicedevoir}[2][Poser et effectuer les calculs suivants]
\begin{multicols}{2}
\begin{enumerate}[itemsep=2em]
\item  $14{,}56+9{,}785$ \\ \encart{5cm}
\item  $656{,}8-33{,}64$ \\ \encart{5cm}
\end{enumerate}
\end{multicols}
\end{exercicedevoir}


\begin{exercicedevoir}[2][Résoudre le problème suivant en détaillant les calculs]
\begin{multicols}{2}
Chez le boulanger, je dois payer  $6{,}40$\,€. \\
Je donne un billet de $10$ \,€. \\Combien me rend-on ? \\
\newcolumn
\encart{5cm}
\end{multicols}
\end{exercicedevoir}


\begin{exercicedevoir}[4][Résoudre le problème suivant en détaillant les calculs]
\begin{multicols}{2}
Un électricien dispose d’un rouleau de fil électrique de $70\text{ m}$. \\ \\
Il découpe $3$ morceaux de fil de ce rouleau de $8{,}50\text{ m}$ chacun.\\ \\
Quelle longueur de fil électrique reste-t-il dans le rouleau ? \\
\newcolumn
\encart{5cm}
\end{multicols}
\end{exercicedevoir}


\begin{exercicedevoir}[4][Résoudre le problème suivant en détaillant les calculs]

\begin{multicols}{2}
Nadia repère des serviettes dans un magazine de publicité à $18{,}94$\,\,€~l'unité. \\ \\
Si Nadia achetait une serviette à $18{,}94$\,\,€~l'unité puis d'autres articles pour $12{,}48$\,\,€, quel serait le prix final\,?  \\
\newcolumn
\encart{5cm}
\end{multicols}
\end{exercicedevoir}


\begin{exercicedevoir}[8][Répondre aux questions suivantes]
Zoé repère dans un magasin une carte de collection à $12{,}15$,€ et un porte-clés à $13{,}17$,€.
Elle souhaite aussi acheter un carnet à $8{,}26$,€.

\medskip

Zoé possède un bon de réduction de $7{,}48$,€ valable sur l'ensemble de ses achats.

\medskip

Quelle somme Zoé devra-t-elle payer au final après la réduction ? 
\end{exercicedevoir}
\vspace{-.5cm}
\encart{5cm}

\end{document}

%%% Local Variables:
%%% mode: LaTeX
%%% TeX-master: t
%%% TeX-master: t
%%% End:

